% Imporditud: tools/import_eksperimendid.py
% Allikas: Eesti füüsikaolümpiaadi eksperimendi ülesannete
%   kogu (Rael Kalda, 2023), ülesanne Ü67, lahendus L67.
\setAuthor{EFO žürii}
\setRound{piirkonnavoor}
\setYear{2018}
\setNumber{P E1}
\setDifficulty{2}
\setTopic{varia}
\setVahendid{valge paberileht (2 lehte)}
\setPunktid{10}
\setAllikas{Eksperimendi kogu Ü67, lahendus lk 63}
\setLahendusseis{toores_ilma_skeemita}

\prob{Paberi mass}
Leidke paberilehe mass, kui on teada, et lühema külje pikkus a = 210 mm. Paberipaki peale on märgitud 80 g/m$^2$.
\textit{Märkus:} Joonlaua kasutamine ei ole lubatud!

\hint

\solu
Kiri ''80 g/m$^2$'' tähendab, et paberi iga ruutmeetri mass on 80 g. Tähistame selle tähega r = 80 g/m$^2$. Kasutades ühikut, saab tuletada valemi r = m/S, kus m on paberilehe mass ja S on pindala. Paberilehe pindala S = lpikemllühem, kus lpikem ja llühem on paberilehe mõõtmed. Paberilehe mass m = rS. Pikema külje leidmiseks voldime ühe paberi külje vähemalt 4 korda kokku (paberi külg jaguneb 16 lõiguks) ning vaatame mitu sellist lõiku mahub kummagi külje sisse. Teades, mitu korda erinevad paberi külgede pikkused, saame leida pikema külje pikkuse.

\textit{Žürii hindamisskeem on kogumikus lk 63; siia ei ole seda ümber kirjutatud.}
\probend
